\chapter{SIMD、指令级并行与向量化}

SIMD 让一条指令处理多个数据 lane，是现代 CPU 单核吞吐的重要来源。指令级并行则让多个互不依赖的操作同时在不同执行资源上推进。二者共同决定很多数值循环的上限。

\section{从标量语义到向量语义}

向量化不是把代码表面改成更复杂的形式，而是确认多个元素之间没有依赖，可以被同一条向量指令处理。数组逐元素加法容易向量化，因为每个元素独立；前缀和不容易直接向量化，因为后一个结果依赖前一个结果。

编译器自动向量化需要知道别名关系、对齐、循环边界和数据依赖。使用 \code{std::span} 表达连续区间有助于接口清晰，但编译器仍可能因为别名保守而放弃向量化。查看优化报告和汇编，比猜测更可靠。

\section{别名和 restrict 思维}

两个指针可能指向同一片内存时，编译器必须保守。例如 \code{a[i] = b[i] + c[i]} 中，如果 a、b、c 可能重叠，编译器要保证重叠情况下结果仍正确，这可能阻碍向量化。C++ 没有标准 \code{restrict} 关键字，但接口设计可以减少别名不确定性，例如使用不重叠容器、清晰文档和局部拷贝。

手写 SIMD 前，应该先让标量代码容易被编译器理解。清晰的循环边界、连续数据、少分支、少别名，是自动向量化的基本条件。

\section{多累加器和延迟隐藏}

即使不用 SIMD，多个独立累加器也能提高吞吐。单累加器形成长依赖链，每次加法依赖上一次结果；多个累加器把依赖链拆开，让乱序执行有更多可调度操作。SIMD 内核中也常见多个向量累加器，目的相同。

这就是并行归约在单线程内部也有优化空间的原因。先分成多个局部和，再合并，不仅适用于多线程，也适用于单核流水线。

\begin{lstlisting}[language=C++,caption={多个累加器拆开依赖链}]
std::int64_t s0 = 0;
std::int64_t s1 = 0;
for (std::size_t i = 0; i + 1 < values.size(); i += 2) {
    s0 += values[i];
    s1 += values[i + 1];
}
const std::int64_t sum = s0 + s1;
\end{lstlisting}

这段代码的意义不是只用两个累加器，而是展示依赖链拆分。真实内核会根据执行端口、寄存器数量和向量宽度选择更多累加器。累加器太少隐藏不了延迟，太多会增加寄存器压力。

\section{尾部处理和对齐}

向量宽度通常不能整除数组长度，因此要处理尾部元素。常见方法是标量尾循环、掩码加载和补齐填充。对齐可以减少某些加载成本，但现代 CPU 对非对齐连续访问已经比较友好；真正要避免的是跨 cache line、跨页和随机访问。

\section{SIMD 与线程的关系}

SIMD 和多线程不是替代关系。SIMD 提高单核吞吐，多线程使用多个核心。高性能程序通常先让单线程内核足够高效，再扩展到多线程。否则多线程只会复制低效内核，并更快撞上内存带宽。

\section{数值语义}

向量化可能改变浮点运算顺序，从而改变舍入误差。整数加法在不溢出的数学模型下更容易重排，浮点加法不严格满足结合律。编译器的 fast-math 选项会放宽数值语义，可能提高性能，也可能改变边界行为。工程中必须明确是否允许这种变化。

\begin{labbox}
实验：用编译器优化报告观察一个数组加法循环是否向量化。再写单累加器和四累加器求和，比较性能。记录编译选项和 CPU 指令集。
\end{labbox}
